% =====================================================================
%  Comandos e ambientes próprios do MOP
%  (usados em todas as atividades -- manter a mesma forma de uso)
% =====================================================================

% ---------------------------------------------------------------------
% \pendencia{texto}
%   Registra, no ponto do texto, uma inconsistência ou lacuna a
%   resolver. Todas aparecem listadas no Apêndice A. Para ocultá-las
%   na versão final, usar \mostrarpendenciasfalse em main.tex.
% ---------------------------------------------------------------------
\newcommand{\pendencia}[1]{%
  \todo[inline, color=laranjaclaro, bordercolor=orange,
        caption={#1}]{\textbf{Pendência:} #1}}

% Lista das pendências (sem título próprio), usada no Apêndice A
\makeatletter
\newcommand{\listapendencias}{\@starttoc{tdo}}
\makeatother

% ---------------------------------------------------------------------
% \fluxograma{arquivo}
%   Insere um fluxograma da pasta figuras/, reduzindo-o apenas se for
%   mais largo que a mancha de texto.
% ---------------------------------------------------------------------
\newcommand{\fluxograma}[1]{\adjustbox{max width=\linewidth}{\input{figuras/#1}}}

% ---------------------------------------------------------------------
% \atividade{3.2.2.x}
%   Grafia padronizada do código de uma atividade no texto corrido.
% ---------------------------------------------------------------------
\newcommand{\atividade}[1]{Atividade~#1}

% ---------------------------------------------------------------------
% Cabeçalho de quadros/tabelas (linha colorida em negrito)
% ---------------------------------------------------------------------
\newcommand{\cab}[1]{\cellcolor{azulclaro}\textbf{#1}}

% Colunas utilitárias para tabularx
\newcolumntype{L}{>{\raggedright\arraybackslash}X}
\newcolumntype{C}{>{\centering\arraybackslash}X}
\newcolumntype{P}[1]{>{\raggedright\arraybackslash}p{#1}}
\newcolumntype{Q}[1]{>{\centering\arraybackslash}p{#1}}

% ---------------------------------------------------------------------
% Estilos TikZ dos fluxogramas
% ---------------------------------------------------------------------
\tikzset{
  etapa/.style={draw=azulpsh, fill=azulclaro, rounded corners=2pt,
                align=center, font=\footnotesize, inner sep=4pt,
                minimum height=0.9cm, text width=#1},
  etapa/.default={3.4cm},
  marco/.style={draw=azulpsh, fill=azulpsh, text=white, rounded corners=2pt,
                align=center, font=\footnotesize\bfseries, inner sep=4pt,
                minimum height=0.9cm, text width=#1},
  marco/.default={7.4cm},
  externo/.style={draw=gray, fill=cinzaclaro, rounded corners=2pt,
                  align=center, font=\footnotesize, inner sep=4pt,
                  minimum height=0.9cm, text width=#1},
  externo/.default={3.4cm},
  decisao/.style={draw=orange!80!black, fill=laranjaclaro, diamond, aspect=2.2,
                  align=center, font=\footnotesize, inner sep=1pt},
  correcao/.style={draw=orange!80!black, fill=laranjaclaro, rounded corners=2pt,
                   align=center, font=\footnotesize, inner sep=4pt,
                   minimum height=0.9cm, text width=#1},
  correcao/.default={2.6cm},
  transversal/.style={draw=verdeclaro!50!black, fill=verdeclaro, dashed,
                      rounded corners=2pt, align=center, font=\footnotesize,
                      inner sep=4pt, text width=#1},
  seta/.style={-{Stealth[length=2.2mm]}, thick, draw=azulpsh!80!black},
  setatracejada/.style={-{Stealth[length=2.2mm]}, thick, dashed, draw=gray},
  rotulo/.style={font=\scriptsize\itshape, fill=white, inner sep=1pt},
  % etapa em forma de seta (linhas do tempo)
  fase/.style={signal, signal to=east, signal from=west, draw=azulpsh,
               fill=azulclaro, align=center, font=\scriptsize\bfseries,
               minimum height=1.1cm, text width=#1, inner sep=2pt},
  fase/.default={2.1cm},
  % caixa de produto ou documento resultante de uma etapa
  produto/.style={draw=azulpsh!70, fill=white, align=center, font=\scriptsize,
                  text width=#1, inner sep=3pt, minimum height=0.7cm,
                  rounded corners=1pt},
  produto/.default={2.1cm},
  % faixa (raia) de responsável em fluxogramas do tipo swimlane
  raia/.style={draw=gray!60, fill=cinzaclaro!70, minimum width=#1,
               minimum height=0.9cm, anchor=north west},
  titulo raia/.style={font=\scriptsize\bfseries, text=azulpsh, align=center,
                      text width=1.6cm}
}
